%%-----------Chapter 2------------------------------------------
\chapter{Dimensional Analysis}

Here is a challenge that sounds impossible. Take a pendulum --- a stone on a
string. Its period (the time for one full swing) surely depends on things like
the length of the string, the mass of the stone, how hard gravity pulls. Now
the challenge: \emph{predict the formula for the period without using any law
of motion at all}. No forces, no Newton, no experiment. Just bookkeeping.

By the end of this short chapter you will do exactly that --- and you will
discover along the way that the mass of the stone \emph{cannot appear in the
answer}, a fact you can then walk outside and test with two stones and a
shoelace. The trick is called \textbf{dimensional analysis}, and physicists
reach for it constantly: to catch algebra mistakes, to guess formulas before
deriving them, and sometimes (as the supplementary box at the end of the
chapter shows) to uncover military secrets from a published photograph.

\begin{predictbox}[2.1: The pendulum, before any theory]
A pendulum of length $l$ carries a bob of mass $m$, swinging under gravity
$g$. Before reading on, write down your instincts:
\begin{enumerate}[leftmargin=2em, itemsep=1pt]
  \item If you double the mass of the bob, does the period get longer,
        shorter, or stay the same?
  \item If you double the length of the string, what happens to the period ---
        doubles? quadruples? something else?
\end{enumerate}
Keep your answers. The chapter will grade them --- twice: once by bookkeeping,
once (if your class has a shoelace) by experiment.
\end{predictbox}

\section{Dimension}

In physics, we measure many different things, but they are all built from a small set of seven
fundamental base quantities (see Table~\ref{tab:base_quantities}).

\begin{table}[H]
  \centering
  \caption{The seven fundamental base quantities, their dimensions, and sample units.}
  \label{tab:base_quantities}
  \begin{tabular}{lcc}
    \toprule
    \textbf{Base Quantity} & \textbf{Dimension} & \textbf{Sample Units} \\
    \midrule
    Mass                          & $[M]$     & kg, g, lbs \\
    Length                        & $[L]$     & m, cm, feet, miles \\
    Time                          & $[T]$     & s, minutes, hours, years \\
    Electric Current              & $[I]$     & A, mA \\
    Thermodynamic Temperature     & $[\Theta]$& K, $^\circ$C, $^\circ$F \\
    Amount of Substance           & $[N]$     & mol \\
    Luminous Intensity            & $[J]$     & cd \\
    \bottomrule
  \end{tabular}
\end{table}

The \textbf{dimension} of any physical quantity represents the powers to which these base
quantities must be raised to express that quantity. Example:
\[
  [\text{Force}] = [\text{mass}]\times\frac{[\text{velocity}]}{[\text{time}]}
                 = [\text{mass}]\times\frac{[\text{length/time}]}{[\text{time}]}
                 = MLT^{-2}.
\]
Such an expression for a physical quantity in terms of the base quantities is called the
\textbf{dimensional formula}. Thus, the dimensional formula of force is $MLT^{-2}$. In other
words, we say the dimensions of force are 1 in mass, 1 in length, and $-2$ in time. The
dimensions in all other base quantities are zero.

Note that in this type of calculation the magnitudes are not considered. It is equality of the
\emph{type} of quantity that enters. Thus, change in velocity, initial velocity, average velocity,
and final velocity are all equivalent in this discussion --- each one is $LT^{-1}$.

The physical quantity that is expressed in terms of the base quantities is enclosed in square
brackets to remind us that the equation is among the dimensions and not among the magnitudes.


\section{The Principle of Homogeneity}

You cannot add apples to oranges, and in physics, you cannot add velocity to force.

\begin{keybox}[The Principle of Homogeneity]
The dimensions of all separated terms in a physical equation must be identical.
\end{keybox}

Let us check the kinematic equation $x = ut + \tfrac{1}{2}at^2$:
\begin{itemize}
  \item $[x] = L$
  \item $[ut] = (LT^{-1})(T) = L$
  \item $[\tfrac{1}{2}at^2] = (LT^{-2})(T^2) = L$
\end{itemize}
Since every term evaluates to length $[L]$, the equation is dimensionally correct. Note: the
constant $\tfrac{1}{2}$ has no dimension and is ignored in this check.

\begin{checkpoint}[2.1: Spot the impostor]
A student solving a problem about a falling stone produces four candidate
formulas, where $v$ is a speed, $g$ an acceleration, $h$ a height, and $t$ a
time. Exactly one of them \emph{could} be right. Which one?
\begin{enumerate}[label=(\alph*), leftmargin=2.5em, itemsep=0pt]
  \item $v = 2gh$
  \item $v = \sqrt{2gh}$
  \item $h = \tfrac12 g t$
  \item $t = \sqrt{2h g}$
\end{enumerate}
\checkanswer{Answer: (b). Check each: $[gh]=(LT^{-2})(L)=L^2T^{-2}$, which is
speed \emph{squared} --- so (a) fails but its square root (b) has dimension
$LT^{-1}$, a speed. In (c), $[gt]=LT^{-1}$ is a speed, not a height. In (d),
$[\sqrt{hg}]=\sqrt{L^2T^{-2}}=LT^{-1}$, again a speed, not a time. Note what
the check cannot do: it cannot tell you whether the 2 in (b) is right.}
\end{checkpoint}


\section{Uses of Dimensional Analysis}

\subsection{Checking the Correctness of an Equation}

As shown above, if an equation is not dimensionally homogeneous, it is physically impossible
and strictly wrong.

\textbf{Limitation:} A dimensionally correct equation might still be wrong if numerical constants
are missing.

\subsection{Conversion of Units}

Dimensions help us find conversion factors between different unit systems (like SI and CGS).
For example, to convert 1 Pascal (SI unit of pressure) to CGS pressure:
\[
  [\text{Pressure}] = \frac{[\text{Force}]}{[\text{Area}]} = \frac{MLT^{-2}}{L^2} = ML^{-1}T^{-2}.
\]
\[
  1\;\text{Pascal} = (1\;\text{kg})(1\;\text{m})^{-1}(1\;\text{s})^{-2}, \qquad
  1\;\text{CGS pressure} = (1\;\text{g})(1\;\text{cm})^{-1}(1\;\text{s})^{-2}.
\]
\[
  \frac{1\;\text{Pascal}}{1\;\text{CGS pressure}}
  = \left(\frac{1\;\text{kg}}{1\;\text{g}}\right)
    \left(\frac{1\;\text{m}}{1\;\text{cm}}\right)^{-1}
    \left(\frac{1\;\text{s}}{1\;\text{s}}\right)^{-2}
  = (10^3)(10^2)^{-1} = 10.
\]
So $1\;\text{Pascal} = 10\;\text{CGS pressure}$.

\subsection{Deducing Relations Among Physical Quantities}

This is the payoff promised at the start of the chapter, and it has a
history. Lord Rayleigh --- the physicist who explained why the sky is blue ---
used this ``method of dimensions'' so often and so effectively in the late
1800s that it became a standard weapon in every physicist's toolkit; he was
still writing notes to \emph{Nature} in 1915 urging that the method could
save ``much labour'' by anticipating results before any hard analysis. The
idea: if we know \emph{which} variables affect a physical quantity, the
requirement that the dimensions balance often pins down \emph{how} they can
combine, up to a dimensionless constant.

Time to grade your pendulum predictions. Suppose the time period $t$ of a simple pendulum
depends on its length $l$, the mass of the bob $m$, and the acceleration due to gravity $g$.
We assume a product relation of the form
\begin{equation}
  t = k\,l^a\,m^b\,g^c,
\end{equation}
where $k$ is a dimensionless constant and $a$, $b$, $c$ are exponents we want to evaluate.
Taking dimensions on both sides:
\[
  T = L^a\,M^b\,(LT^{-2})^c = L^{a+c}\,M^b\,T^{-2c}.
\]
Equating powers on both sides:
\begin{itemize}
  \item For $M$: $b=0$
  \item For $T$: $-2c=1 \implies c=-\tfrac{1}{2}$
  \item For $L$: $a+c=0 \implies a=\tfrac{1}{2}$
\end{itemize}
Substituting these values gives
\begin{equation}
  t = k\sqrt{\frac{l}{g}}.
\end{equation}
In the actual formula $k=2\pi$, which is determined from the full physical analysis, not from
dimensional analysis alone.

Look at what the bookkeeping just told us. The mass exponent came out
$b=0$: \emph{the mass of the bob cannot appear in the formula}. A heavy stone
and a light stone on equal strings must swing in step --- not because we
measured it, but because no other answer can balance dimensionally. (If your
first instinct in Stop and Predict 2.1 was ``heavier swings slower,'' you are
in good company --- most people guess that. Now tie two different stones to
equal shoelaces and watch.) And doubling the length does not double the
period; it multiplies it by $\sqrt{2}\approx 1.4$, because the length enters
under a square root. This ``predict by bookkeeping, then test'' cycle is
exactly how you will verify the formula in the lab later this year.

\textbf{Limitations:}
\begin{enumerate}
  \item It only works for \emph{product-type} dependencies. It cannot deduce equations where
        terms are added or subtracted (like $x = ut + \tfrac{1}{2}at^2$).
  \item It cannot find numerical constants. The dimensional method left us with an
        unknown constant $k$ in the pendulum equation. We need actual experiments (or deeper
        theory) to find that $k=2\pi$.
\end{enumerate}

\begin{supplementary}{How dimensional analysis leaked an atomic secret}
In 1945 the United States exploded the first atomic bomb in the New Mexico
desert. The energy released --- the ``yield'' --- was one of the most closely
guarded secrets in the world. Two years later, the government declassified
something that seemed harmless: a filmstrip of photographs of the growing
fireball, each frame stamped with the time after detonation and a distance
scale.

The British physicist G.\,I.\ Taylor looked at those photographs and reasoned
like this. Once the blast is underway, the radius $R$ of the fireball can
plausibly depend on only three things: the energy $E$ released, the density
$\rho$ of the surrounding air, and the time $t$ since detonation. There is
exactly one way to combine $[E]=ML^2T^{-2}$, $[\rho]=ML^{-3}$, and $[T]$ into
a length (try it --- it is the same game we played with the pendulum):
\[
  R \sim \left(\frac{E\,t^2}{\rho}\right)^{\!1/5}
  \qquad\Longleftrightarrow\qquad
  E \sim \frac{\rho\,R^5}{t^2}.
\]
Taylor read $R$ and $t$ off the published frames (about 140\,m at
$t=25$\,ms), put in $\rho\approx1.2\;\mathrm{kg/m^3}$ for air, and out came
$E\sim10^{14}$\,J --- the equivalent of roughly twenty thousand tonnes of TNT.
When he published the estimate in 1950 it was embarrassingly close to the
still-classified value. The American military was reportedly not amused; no
law had been broken, only bookkeeping applied. \emph{Wonder en is gheen
wonder.}
\end{supplementary}
